% ==============================================================================
% STEM (Science / Technology / Engineering / Mathematics) Paper Template
% Structure: IMRaD  (Introduction / Methods / Results and Discussion)
% Standardized by ANSI Z39.16-1972; introduction follows Swales' CARS model
% (Create A Research Space, Swales 1990).
%
% This skeleton uses the standard `article` class + IEEEtran.bst so it compiles
% on any TeX Live / MiKTeX install.  For journal-specific submission, swap the
% preamble with the matching journal template under ../journals/ .
%
% Compile:  pdflatex main -> bibtex main -> pdflatex main -> pdflatex main
% ==============================================================================
\documentclass[11pt,a4paper]{article}

% ---- Page geometry ----
\usepackage[a4paper,top=25mm,bottom=25mm,left=25mm,right=25mm]{geometry}

% ---- Fonts & encoding ----
\usepackage[T1]{fontenc}
\usepackage[utf8]{inputenc}
\usepackage{mathptmx}             % Times for body + math
\usepackage{helvet}
\usepackage{courier}

% ---- Math, figures, tables ----
\usepackage{amsmath,amssymb,amsfonts}
\usepackage{amsthm}               % theorem/lemma/proof environments
\usepackage{graphicx}
\usepackage{booktabs}
\usepackage{array}
\usepackage{xcolor}
\usepackage{algorithm}            % algorithm floating environment
\usepackage{algorithmic}          % IEEE-style pseudo-code commands

% ---- Theorem-like environments (optional, for theory papers) ----
\newtheorem{theorem}{Theorem}
\newtheorem{lemma}[theorem]{Lemma}
\newtheorem{proposition}[theorem]{Proposition}
\newtheorem{corollary}[theorem]{Corollary}

% ---- Citations: numeric [1], [2], [3]-[5] in citation order ----
\usepackage[numbers,sort&compress]{natbib}
\bibliographystyle{IEEEtran}      % swap with elsarticle-num / ACM-Reference-Format / siamplain / unsrtnat as needed

% ---- Hyperref last ----
% \usepackage{hyperref}
% \hypersetup{colorlinks=true,linkcolor=blue,citecolor=blue,urlcolor=blue}

% ---- Line numbers for review (uncomment for double-blind submission) ----
% \usepackage{lineno} \linenumbers

\title{A Concise and Informative Title of the STEM Submission}
\author{First Author\textsuperscript{1,*}, Second Author\textsuperscript{2}, Senior Author\textsuperscript{1}\\
\textsuperscript{1}Department of X, University of Y, City, Country\\
\textsuperscript{2}Department of Z, Institute of W, City, Country\\
\textsuperscript{*}e-mail: first.author@univ.edu}
\date{}

\begin{document}
\maketitle

% ==============================================================================
% Abstract  (150-250 words; no citations; state problem / approach / results / significance)
% ==============================================================================
\begin{abstract}
% TODO: 150-250 words. Sentence 1-2: problem statement. Sentence 3-5: approach.
% Sentence 6-8: key quantitative results with effect sizes. Sentence 9: significance.
\end{abstract}

\noindent\textbf{Keywords:} % TODO: 4-6 keywords separated by commas.
keyword one, keyword two, keyword three, keyword four

% ==============================================================================
% 1. Introduction  (CARS model: territory -> niche -> occupy niche)
% ==============================================================================
\section{Introduction}
\label{sec:intro}
% TODO intro-CARS Move 1 (territory): claim centrality of the research area
% and review recent prior work with \citep{ref_example}.
% TODO intro-CARS Move 2 (niche): indicate a gap / contradiction / limitation
% in the prior work reviewed above.
% TODO intro-CARS Move 3 (occupy): state the purpose, the bulleted
% contributions, and the paper roadmap. Use \citep{} for parenthetical cites.

% ==============================================================================
% 2. Related Work
% ==============================================================================
\section{Related Work}
\label{sec:related}
% TODO related-work: cluster prior work by approach family; for each cluster
% compare methodology, data, and performance; explicitly position this work.

% ==============================================================================
% 3. Methods / System Model
% ==============================================================================
\section{Methods}
\label{sec:methods}

\subsection{Problem formulation}
% TODO: state the problem formally; introduce notation; list assumptions.

\subsection{Proposed method}
% TODO: describe the proposed method step by step; reference Algorithm 1.

\begin{algorithm}[t]
\caption{Outline of the proposed method}
\label{alg:method}
\begin{algorithmic}[1]
\REQUIRE dataset $\mathcal{D}$, hyper-parameters $\Theta$
\STATE initialise model parameters $\theta$
\WHILE{not converged}
    \STATE $\theta \leftarrow \theta - \eta \,\nabla \mathcal{L}(\theta;\mathcal{D})$
\ENDWHILE
\RETURN trained model with parameters $\theta$
\end{algorithmic}
\end{algorithm}

\subsection{Complexity analysis}
% TODO: time complexity and space complexity in big-O notation.

% ==============================================================================
% 4. Experiments / Results
% ==============================================================================
\section{Experiments}
\label{sec:experiments}

\subsection{Experimental setup}
% TODO experiment-setup: datasets (with citations), baselines, evaluation
% metrics, hardware (CPU/GPU/RAM), software versions, random seed.

\subsection{Main results}
% TODO results: present quantitative results; reference Table 1 and Figure 1.
% Report mean $\pm$ std over multiple runs; state statistical tests.

\begin{table}[t]
\caption{Main results. Caption ABOVE the table (STEM convention).}
\label{tab:results}
\centering
\begin{tabular}{lcc}
\toprule
Method   & Metric A       & Metric B       \\
\midrule
Baseline & 0.812          & 0.754          \\
Ours     & \textbf{0.891} & \textbf{0.832} \\
\bottomrule
\end{tabular}
\end{table}

\begin{figure}[t]
\centering
% \includegraphics[width=0.9\columnwidth]{figures/result.pdf}
\fbox{\rule{0pt}{4cm}\rule{0.9\columnwidth}{0pt}}
\caption{Single-column figure. Caption goes BELOW the figure. Use
\texttt{\textbackslash columnwidth} or \texttt{\textbackslash textwidth}.}
\label{fig:result}
\end{figure}

\subsection{Ablation study}
% TODO: remove or replace each component to verify its individual contribution.

\subsection{Statistical significance}
% TODO: report t-test / Wilcoxon / bootstrap 95\% CI; state $p$-values.

% ==============================================================================
% 5. Discussion
% ==============================================================================
\section{Discussion}
\label{sec:discussion}
% TODO: interpret results, reconcile with prior work, state limitations
% (TODO: limitations) and threats to validity.

% ==============================================================================
% 6. Reproducibility Statement  (STEM-specific)
% ==============================================================================
\section*{Reproducibility Statement}
% TODO reproducibility: state where code, data, pre-trained models, and
% configuration files are deposited (GitHub URL + Zenodo DOI). Disclose the
% random seed, framework versions, and hardware. If any data is restricted,
% state the access procedure.

% ==============================================================================
% 7. Conclusion
% ==============================================================================
\section{Conclusion}
\label{sec:conclusion}
% TODO: recap contributions and the headline result; one sentence of future work.

% ==============================================================================
% Acknowledgements
% ==============================================================================
\section*{Acknowledgements}
% TODO: funding agencies, grant numbers, individuals.

% ==============================================================================
% References  —  IEEEtran numeric style [1], [2], [3]-[5]
% Compile:  pdflatex main -> bibtex main -> pdflatex main -> pdflatex main
% ==============================================================================
\bibliography{references}
% Inline fallback (uncomment if not using BibTeX):
% \begin{thebibliography}{99}
% \bibitem{ref_example} F.~Author and S.~Author, ``Title of the paper,''
%   \emph{Journal Name} XX (20YY) 1--10.
% \end{thebibliography}

\end{document}
